\chapter{Background}
\label{background}

Football is the most popular sport in the world.
Originating in England during the nineteenth century, the game rapidly spread across continents and became deeply embedded in national cultures.
In Brazil, where the sport arrived in 1894, football transcended its role as mere entertainment to become a defining element of national identity.
It is a passion that unites millions of people regardless of social background, region, or creed.

This popularity has also made football a significant economic force.
Beyond matchday revenues, merchandise sales, and broadcasting rights, the global sports betting market has grown substantially in recent years.
Football accounts for the largest share of this market, which is projected to reach approximately 187 billion dollars by 2030 \cite{gvr_2024}.
Questions such as ``which team is better?'' or ``is player A superior to player B?'' are no longer confined to casual conversations among fans.
They now carry financial implications for betting operators and bettors alike, creating demand for rigorous quantitative analysis.

This context has motivated the development of probabilistic models capable of estimating team strengths, predicting match outcomes, and quantifying the uncertainty associated with these predictions.
Such models provide a framework for transforming historical match data into actionable insights, whether for sports analysts, club managers, or betting markets.
This chapter reviews the main modelling approaches proposed in the literature, covering paired comparison models, goal-based formulations, player-level extensions, and a simulation-based validation technique.

Among paired comparison models, the Bradley-Terry model \cite{bradley_terry_1952} was originally developed for analysing experiments where items are compared in pairs and a preference or outcome is recorded.
In the context of sports, this model provides a natural framework for modelling match outcomes by relating the probability of one team beating another to latent skill parameters.
The original formulation considered binary outcomes, representing a win for one of the two items being compared.
\textcite{davidson_1970} extended the Bradley-Terry model, introducing a tie parameter.
This parameter allows the model to assign non-zero probability to draw outcomes, making it more suitable for football applications.
\textcite{davidson_1977} further extended the model to incorporate order effects, which can capture home effects in sports contexts.
These extensions transformed the Bradley-Terry model from a simple binary comparison tool into a flexible framework for modelling sports competitions.
More recently, \textcite{cattelan_2013} developed dynamic Bradley-Terry models that allow team strengths to evolve over time, capturing the temporal dynamics of team performance throughout a season.
This work demonstrated the flexibility of the Bradley-Terry framework for modelling sports tournaments and showed how temporal variation can be incorporated while maintaining the model's interpretability.

Rather than modelling only match outcomes, an alternative and complementary approach uses the number of goals scored by each team as variables to be modelled.
This approach was formalised by \textcite{maher_1982}, where these variables are modelled as Poisson with rates that depend on team offensive and defensive abilities.
We call these Poisson-based formulations as Poisson models, where the focus is modelling explicitly the final score of a match.
\textcite{lee_1997} applied Poisson models to English Premier League data, with a particular emphasis on incorporating home effects as a distinct parameter in the model.
Independently, \textcite{dixon_1997} extended this class of models by introducing a correlation parameter to account for the dependence between home and away goal counts.
Their model incorporated a bivariate Poisson distribution to capture this correlation, building upon earlier work on bivariate Poisson distributions \cite{holgate_1964}.

\textcite{karlis_2003} further developed the bivariate Poisson approach, demonstrating that accounting for correlation improves model fit and predictive performance.
\textcite{karlis_2009} proposed using the Skellam distribution to model goal differences directly, offering an alternative to bivariate Poisson models that focuses on the difference between goal counts rather than the counts themselves.
This approach can be particularly useful when the research question focuses on match outcomes rather than exact scorelines.
Within a Bayesian framework, \textcite{baio_2010} proposed a hierarchical model for estimating team strengths and predicting match scores, introducing a mixture formulation to address the overshrinkage produced by standard hierarchical priors.
Similarly to the dynamic extensions in the Bradley-Terry framework, \textcite{rue_2000} introduced a Bayesian Poisson model with time-varying team strengths, and \textcite{koopman_2015} combined the bivariate Poisson formulation with state-space methods to model the evolution of team abilities across a season.
Recent applications include Poisson-based analyses of the English Premier League \cite{nguyen_2021} and double Poisson formulations for international tournaments such as Euro 2020 \cite{penn_2022}.

Most of the literature on football modelling operates at the club level, treating each match as a confrontation between two atomic entities, where each team possesses an estimated strength parameter derived from observed data.
However, since clubs are formed by a collection of players, it is natural to consider individual-level modelling, where each player is treated as an atomic entity for the model.
In this direction, \textcite{wolf_2020} proposed a player rating system based on the Elo algorithm, originally developed by \textcite{elo_1961} and still used today by the Fédération Internationale des Échecs (FIDE) for chess rankings.
Since the Elo update rule corresponds to stochastic gradient ascent on the Bradley-Terry log-likelihood \cite{aldous_2017}, the work of \textcite{wolf_2020} can be viewed as a player-level extension of the Bradley-Terry framework.
Using data from 18 European competitions from seasons 2014/15 to 2017/18, they calculated player rankings and generated insights into each player's impact on their club's performance.

A critical aspect of statistical modelling is ensuring that the methods used for posterior inference correctly reproduce the intended posterior distributions.
When this is not the case, the resulting parameter estimates and uncertainty quantification can be misleading, potentially leading to incorrect conclusions.
Traditional diagnostics, such as trace plots and effective sample size, can detect some problems but may miss systematic biases.
Nevertheless, they provide no direct guarantee that the samples are drawn from the correct posterior distribution.

\textcite{cook_2006} introduced the concept of Simulation-Based Calibration (SBC) as a method for validating Bayesian inference procedures.
The key insight is that if an inference procedure correctly reproduces the posterior distribution, then the rank of the true parameter value among posterior samples should follow a uniform distribution.
This property provides a powerful diagnostic tool: systematic deviations from uniformity indicate that the inference procedure is not correctly reproducing the posterior distribution.
\textcite{talts_2020} formalized this approach and provided a comprehensive framework for applying SBC to validate Bayesian models.
They showed that SBC can detect various types of problems, including bias, under-dispersion, and overdispersion in posterior distributions.
More recently, \textcite{modrak_2025} extended SBC to high-dimensional settings, by testing uniformity of test quantities.

In the context of Brazilian football, some initiatives include the \textit{Chance de Gol} website \cite{chancedegol}, which has been providing match outcome probabilities since 1999, and the \textit{Probabilidades no Futebol} project \cite{ufmg_futebol}, maintained by the Department of Mathematics at the Universidade Federal de Minas Gerais (UFMG) since 2005.
On the academic side, \textcite{diniz_2018} compared several probabilistic predictive models, including Bayesian multinomial-Dirichlet formulations, evaluating their predictive performance on \textit{Campeonato Brasileiro Série A} (also known as \textit{Brasileirão}) matches using proper scoring rules and calibration diagnostics.

The present work makes three main contributions to the literature on Brazilian football modelling.
First, a comprehensive dataset of Brazilian football matches spanning multiple seasons and competitions is introduced.
Second, Simulation-Based Calibration is applied to validate Bayesian inference in football models.
Finally, club-level formulations are extended to the player level, enabling a more granular analysis of team performance.
